% !TeX root = ../../tesis.tex
%%%%%%%%%%%%%%%%%%%%%%%%%%%%%%%%%%%%%%%%%%%%%%%%%%%%%%%%%%%%%%%%%%%
% §6.3.1 — LOS ESCENARIOS DEL API-I
%
% Estructura:
%   · Entrada: por qué dos casos (mismo trazo, dos modos de operación)
%   · Caso 2: Anillo MB — razones (las da el autor), vía confinada,
%     mapa Propuesta_Anillar_Extendido_MB
%   · Caso 3: Anillo Metro — razones (las da el autor), vía exclusiva,
%     mapa Propuesta_Anillar_Extendido_Metro (trasladado desde §6.1.3)
%   · Cierre: cuadro comparativo de los casos en términos urbanos,
%     sin adelantar resultados
% Las notas de cuadrante de los mapas extendidos traen resultados
% (B(v), T): no usarlos aquí; van en §6.3.2–§6.3.5.
%%%%%%%%%%%%%%%%%%%%%%%%%%%%%%%%%%%%%%%%%%%%%%%%%%%%%%%%%%%%%%%%%%%

% Figura trasladada desde §6.1.3 (2026-10-18); se integra al
% redactar el Caso 3.
\begin{figure}[H]
    \centering
    \includegraphics[width=0.92\textwidth]{%
        Figures/Mapas/Qgis_Resources/RedConAnillo/Escenario_Metro/Base/Propuesta_Anillar_Extendido_Metro}
    \caption[Cuadrantes del corredor anillar]{%
    Corredor anillar propuesto por cuadrante (norponiente, nororiente,
    surponiente y suroriente) y sistemas de transporte que cruza.
    Corrida: 2026-10-04.}
    \label{fig:propuesta-anillo-cuadrantes}
    \fuentefigura{Elaboración propia con \textbf{QGIS};
    datos: \texttt{Apímetro} y \texttt{VFTModel}~\citep{cabrera2025apimetro, cabrera2026VFTModel}.}
\end{figure}
